\documentclass[11.5pt]{article}
\usepackage[letterpaper,margin=22mm]{geometry}
\usepackage[T1]{fontenc}
\usepackage[utf8]{inputenc}
\usepackage{lmodern}
\usepackage{amsmath,amssymb,amsfonts,mathtools}
\usepackage{booktabs,longtable,array,tabularx}
\usepackage{graphicx}
\usepackage{caption}
\captionsetup{font=small,labelfont=bf,labelsep=period}
\usepackage{xcolor}
\usepackage{enumitem}
\usepackage[hidelinks]{hyperref}
\usepackage{fancyhdr}
\usepackage[normalem]{ulem}
\setlength{\parindent}{1.5em}\setlength{\parskip}{0pt}
\definecolor{accent}{HTML}{5B4BD6}
\pagestyle{fancy}\fancyhf{}
\fancyhead[L]{\small\color{gray}The Mechanics and Environmental Significance of the Water Cycle}
\fancyfoot[C]{\small\thepage}
\renewcommand{\headrulewidth}{0.4pt}
\setlength{\emergencystretch}{3em}
\begin{document}
\begin{center}
{\LARGE\bfseries\color{accent} The Mechanics and Environmental Significance of the Water Cycle}\\[6pt]
\end{center}
\vspace{6pt}
\section{Introduction and Scope}
\section{Physical Processes of the Hydrologic Cycle}
\subsection{Phase Changes and Atmospheric Fluxes}
Evaporation is the phase change of water from liquid to gas, driven primarily by solar radiation, which facilitates the exchange of approximately 600 calories of energy per gram of water. This process occurs across oceans, lakes, soil, vegetation, and snow. The ocean serves as the source of 86 percent of global evaporation. Condensation involves the transformation of water vapor back into liquid or solid particles, typically forming clouds, fog, or dew when air is cooled to its saturation point. This process releases the latent heat previously absorbed during evaporation. Precipitation describes the fall of water particles from the atmosphere to the ground, occurring through coalescence or ice-crystal formation. Globally, about 505,000 cubic kilometers of water fall as precipitation annually, with 78 percent of this occurring over the oceans, satisfying the relation \[ P_{\text{annual}} \approx 505,000 \text{ km}^3/\text{yr} \quad (\text{Marine Portion} \approx 78\\%) \].

\subsection{Terrestrial and Subsurface Hydrology}
Interception refers to the interruption of water movement by vegetation or surface depressions before it reaches the soil. When vegetation is present, water collects on organic surfaces until its weight exceeds surface tension, causing throughfall. Infiltration is the movement of water across the soil-atmosphere boundary, a process governed by soil porosity, permeability, surface texture, and initial moisture content. Once water enters the soil, percolation facilitates its downward movement through soil layers via gravity and capillary forces, ultimately reaching the water table that separates the vadose zone from the zone of saturation where groundwater is stored in aquifers.

\subsection{Biological and Channel Dynamics}
\section{Global Distribution and Reservoir Residence Times}
\section{Environmental Significance and Climate Interactions}
The global water cycle acts as the primary thermodynamic engine regulating Earth's climate and distributing freshwater resources essential for terrestrial ecosystems. Anthropogenic emissions and climate warming have fundamentally altered these dynamics. As atmospheric temperatures rise, the water cycle accelerates in accordance with the Clausius-Clapeyron relation, which dictates that saturation vapor pressure increases by approximately $7$ percent per degree Celsius of warming. This thermodynamic enhancement results in intensified precipitation events, elevated flood frequencies, and concurrently aggravated drought conditions in continental interiors due to accelerated soil moisture evaporation. Furthermore, thermal expansion of warming ocean waters combined with the melting of land-based glaciers has driven global mean sea levels higher by $0.10$ to $0.20$ meters over the past century, with projected increases reaching up to $1.1$ meters under high-emission scenarios.

\section{References}
National Oceanic and Atmospheric Administration (NOAA), Northwest River Forecast Center. Description of the Hydrologic Cycle. Available at: https://www.nwrfc.noaa.gov/info/water\_cycle/hydrology.cgi.\textbackslash{}n\textbackslash{}nIntergovernmental Panel on Climate Change (IPCC). Sixth Assessment Report: Climate Change 2021: The Physical Science Basis. Cambridge University Press.\textbackslash{}n\textbackslash{}nUniversity Corporation for Atmospheric Research (UCAR). Water Cycle and Climate Change Learning Zone. Available at: https://scied.ucar.edu/learning-zone/climate-change-impacts/water-cycle-climate-change.\textbackslash{}n\textbackslash{}nNASA Global Precipitation Measurement (GPM) Mission. Why Are Water Cycle Processes Important? Available at: https://gpm.nasa.gov/resources/faq/why-are-water-cycle-processes-important.


\begin{thebibliography}{9}
\bibitem{ref45} Description of the Hydrologic Cycle - National Oceanic and Atmospheric Administration. https://www.nwrfc.noaa.gov/info/water\_cycle/hydrology.cgi.
\bibitem{ref46} Water cycle - Wikipedia. https://en.wikipedia.org/wiki/Water\_cycle.
\bibitem{ref47} water cycle climate change. https://scied.ucar.edu/learning-zone/climate-change-impacts/water-cycle-climate-change.
\bibitem{ref48} why are water cycle processes important. https://gpm.nasa.gov/resources/faq/why-are-water-cycle-processes-important.
\end{thebibliography}
\end{document}
